\section{The K230 foundation and its verification boundary}
\label{sec:k230-foundation}

The user-specified K230 repository supplies a substantially smaller
\emph{foundational} starting point: its published machine already has
230 working states, and the theory searched by that machine has a complete
Rocq calibration to ordinary ZFC.  The version inspected here is
\texttt{3142e4dd980f0be12f4abb636c4ae2528afc17eb} of
\url{https://github.com/elliotglazer/k230-rocq}.
The published table's SHA-256 is
\begin{center}\scriptsize\ttfamily
e7ccbcec074534cfd955f0ab94d16f13594e81974074309dc2ce663b9dfd7ba8.
\end{center}
This is an inherited result, not a new construction of the present work.
It gives the appropriate calibration to retain when transferring
compatible source rearrangements or transition-state reductions.

\subsection{One schema, with an essential unconditional closure clause}

Let $\varphi(y,p,\vec q)$ be any first-order formula with output variable
$y$, input variable $p$, and an arbitrary finite tuple of parameters.
Define
\begin{align*}
 p\subseteq u&\quad:\!\!\Longleftrightarrow\quad
       \forall z\,(z\in p\longrightarrow z\in u),\\
 W_\varphi(p,C)&\quad:\!\!\Longleftrightarrow\quad
       (\exists y\,\varphi(y,p,\vec q))
       \longrightarrow\exists y\,(y\in C\land\varphi(y,p,\vec q)),\\
 H_\varphi(C)&\quad:\!\!\Longleftrightarrow\quad
       \forall u\in C\left[
       (\exists p\,\varphi(u,p,\vec q))\land
       \forall p\,(p\subseteq u\longrightarrow W_\varphi(p,C))\right].
\end{align*}
Together with equality axioms, the ordinary theory contains every
universal closure of
\begin{equation}
 \forall a\,\forall x\,\exists C\,
 \bigl[(\varphi(x,a,\vec q)\longrightarrow x\in C)
          \land H_\varphi(C)\bigr].
 \label{eq:k230-ordinary-schema}
\end{equation}
These definitions present \texttt{KPrimeOrd} in ordinary notation.
The formal definitions are in \texttt{KPrimeCore.v} and
\texttt{KPrimeTheory.v}.  In particular, $H_\varphi(C)$ is required
\emph{even when the selected seed instance is false}.  Moving it
inside the seed implication would change the theory.  The actual emitted
formula reorders conjunctions and moves a support witness across a
conjunct in which that witness does not occur; the decoder bridge proves
the correspondence.
\footnote{Pinned source overview:
\url{https://github.com/elliotglazer/k230-rocq/blob/3142e4dd980f0be12f4abb636c4ae2528afc17eb/docs/proof-overview.md}.}

The power of the schema can be seen in short exact instances.  A false
matrix produces an empty raw object, since the support clause excludes
every member of $C$.  For a fixed $a$, the matrix $\varphi(y,p):y=a$,
with seed $x=a$, produces a singleton of $a$.  For raw power set, use
\[
                 \varphi(y,p;a)\quad:\!\!\Longleftrightarrow\quad
                    y=p\land p\subseteq a
\]
with seed input and output both $a$.  The seed gives $a\in C$;
support makes every member of $C$ a raw subset of $a$.  Conversely,
if $p\subseteq a$, closure at the member $a$ requires a witness for
the input $p$, and the only allowed output is $y=p$.  Hence
$p\in C$, giving exactly $C=\{p:p\subseteq a\}$.

After singletons are available, a transitive container of $a$ comes
from the matrix saying that $p$ is a singleton of $y$, seeded by
output $a$ and input a singleton of $a$.  If $v\in u\in C$,
then a singleton of $v$ is a subset of $u$; the closure requirement
therefore puts $v$ into $C$.  These constructions are the formal
patterns in \texttt{KPRaw.v} and \texttt{KPTransitive.v}.
They do not assume that raw objects are extensional sets.
\footnote{The power and singleton derivations are in
\url{https://github.com/elliotglazer/k230-rocq/blob/3142e4dd980f0be12f4abb636c4ae2528afc17eb/coq/settheory/KPRaw.v};
the transitive-container derivation is in the adjacent
\texttt{KPTransitive.v}.}

\subsection{Exact consistency target and the reusable theorem chain}

The formal target is defined literally by
\[
 \mathsf{StandardZFC}(\psi)=
       (\mathsf{ZFc}(\psi)\lor\psi=\mathsf{ax\_reg}')
                       \lor\psi=\mathsf{ax\_ac}.
\]
Thus Foundation and Choice are included.  The development proves
equivalence of its finite-refutation predicate with the corresponding
predicate for both $\mathsf{KPrimeOrd}$ and
$\mathsf{KPrimeOrdIR}$, the latter adding
$\forall x\,\neg(x\in x)$.
Irreflexivity is used for the machine's numerical falsity alias;
its absence from the ordinary lower interpretation does not license
deleting numerical logical rows.

The lower consistency reduction constructs a tagged domain, relation
codes, and a bisimulation quotient satisfying extensional ZF.  Foundation
and Choice are transferred using the pinned upstream relative-consistency
proofs.  The upper reduction interprets each finite schema instance in
the constructible hierarchy, using reflection to obtain the required
closure bounds and Separation to retain supported outputs.
In the quotient power-set step, the proof bounds canonical representatives
before including their equivalent representatives; it does not assume
that every raw code lies in the same naive bound.
\footnote{The actual refutation equivalences are in
\url{https://github.com/elliotglazer/k230-rocq/blob/3142e4dd980f0be12f4abb636c4ae2528afc17eb/coq/settheory/KUpperConsistency.v};
the target and lower bridge are in the adjacent \texttt{KQuotZFC.v}.}

The released final theorem is assembled by the following explicit chain:
\begin{align*}
 \operatorname{Halt}(K230,\mathrm{blank})
 &\ \longleftrightarrow\ \operatorname{RMHalt}(p_{230},R_0)
                      &&\texttt{tm230\_rm},\\
 &\ \longleftrightarrow\ \operatorname{EnumHalt}(\mathsf{seed230})
                      &&\texttt{prog230},\\
 &\ \longleftrightarrow\ \mathsf{thm230}(1)
                      &&\texttt{seed230\_iff\_thm},\\
 &\ \longleftrightarrow\ \mathsf{StandardZFC}\vdash_{\mathrm{TC}}\bot
                      &&\texttt{thm230\_iff\_standardZFC}.
\end{align*}
The seed is
\[
  (\mathit{prooflist},\mathit{nextproof},\mathit{top},\mathit{stack},
     \mathit{parameter},\mathit{cache})=(0,3,0,0,3,0),
\]
not an empty initial search state.  The published blank bootstrap is
408,419 steps.  Its physical register frame has first nine values
$[0,0,0,3,0,0,3,0,0]$, followed by 167 spare values equal to~1.

A new table that provably implements the same $p_{230}$ and seed can
reuse the last three arrows.  A change of source tree needs a complete
control-graph equivalence, including every jump and guarded edge.
An all-tape simulation between whole machines can instead transport the
published final theorem directly.  The 235-state ancestor's CP/GIC
source is a different source program, so its old logical bridge cannot
be silently substituted for the K230 bridge.
\footnote{The final composition is exactly the short proof in
\url{https://github.com/elliotglazer/k230-rocq/blob/3142e4dd980f0be12f4abb636c4ae2528afc17eb/coq/tm/Correct230.v}.}

\subsection{Independent input audit and retained Linux build evidence}

This session independently checked the 103 local manifest entries,
including all 101 local Rocq proof files, and compared all 460 literal
transition rows against \texttt{Table230.v}.  The source and table
hashes match.  After nested-comment removal, no local proof file
contains an axiom declaration, admitted proof, aborted proof, or unchecked
cast token from the audit's listed set.  All seven existing build-script
tests pass; they test infrastructure contracts, not the mathematical
theorems.

A separate literal-source parser also expands all 5,124 rows of
\texttt{Prog230.v} and verifies exact ordered equality of its 5,125
leaves with the published JSON and the normalized working source.
It checks the edge translation, root, node bound, default, and initial
register expression explicitly.  The expansion is acyclic, every
node is reachable, and its maximum leaf depth is 40.  This closes
the concrete JSON-to-formal-program correspondence used by the source
transport argument in Section~\ref{sec:source229-independent}.

More substantially, the GitHub plugin retrieved the successful Linux
workflow for the exact inspected commit, run
\texttt{39258562285}.  The downloaded
\texttt{k230-linux-verification} artifact, identifier
\texttt{12192455062}, has SHA-256
\begin{center}\scriptsize\ttfamily
60702d530e637fd343eda4d370b73d71924ca5c3f581444196acdf5da344f15f.
\end{center}
Its retained records match the inspected manifest, machine, and upstream
commit.  They record Rocq~9.1.0, \texttt{rocq-stdlib}~9.0.0, and
OCaml~4.14.2; compilation of all 216 proof sources; two final
\texttt{Closed under the global context} reports; and a successful
independent \texttt{coqchk} run with exit status zero.  The reported
build-plus-check duration is 1,838.21 seconds, including 516.58 seconds
for the independent kernel checker.
\footnote{Exact successful workflow:
\url{https://github.com/elliotglazer/k230-rocq/actions/runs/39258562285}.
The unmodified downloaded JSON records and logs are retained with this
research delivery.}

The checker log also lists three axioms present in the loaded environment:
\nolinkurl{BusyCoq.LibTactics.inj_pair2},
\texttt{proof\_irrelevance}, and \texttt{eq\_rect\_eq}.
This does not contradict the two closed final dependency reports:
the former list describes the loaded environment, while
\texttt{Print Assumptions} describes the actual dependencies of each
specified theorem.  The final theorems use none of those declarations.
The retained checker summary also reports no type-in-type dependency,
unsafe fixpoint, or assumed positivity.

No local Rocq toolchain was installed at the time of this source audit;
the complete proof rebuild was performed by the identified upstream
Linux workflow, whose artifact was retrieved and checked here.
The endpoint is an external Rocq theorem about the exact K230 table
and a syntactic finite-refutation predicate.  It is not an internally
arithmetized PA derivation, and it does not by itself constitute a new
Rocq endpoint for any reduced table.  Any ordinary mathematical
transport proof for such a table must be stated separately.
